\section*{Capítulo \ref{ch:g10:vectors} --- Vectores}

\begin{solution}{exo:g10:vectors:1}
$\vect{AB} + \vect{BD} = \vect{AD}$ (Chasles).

$\vect{KL} + \vect{LM} + \vect{MK} = \vect{KK} = \vec 0$.

$\vect{AC} - \vect{BC} = \vect{AC} + \vect{CB} = \vect{AB}$.

$\vect{AB} + \vect{CA} = \vect{CA} + \vect{AB} = \vect{CB}$.
\end{solution}

\begin{solution}{exo:g10:vectors:2}
$\vect{AB}\,(5-2,\ 3-1) = (3, 2)$; \quad
$\vect{BC}\,(-1-5,\ 4-3) = (-6, 1)$; \quad
$\vect{AC}\,(-1-2,\ 4-1) = (-3, 3)$.
Comprobación: $(3, 2) + (-6, 1) = (-3, 3)$, coordenada a coordenada.
\end{solution}

\begin{solution}{exo:g10:vectors:3}
$\vec u + \vec v = (2 + (-1),\ -3 + 4) = (1, 1)$.

$3\vec u = (6, -9)$.

$2\vec u - \vec v = (4 - (-1),\ -6 - 4) = (5, -10)$.

$\norm{\vec u} = \sqrt{2^2 + (-3)^2} = \sqrt{13}$.
\end{solution}

\begin{solution}{exo:g10:vectors:4}
(a) $4 \times 9 - 6 \times 6 = 36 - 36 = 0$: \omterm{def:g10:vectors:collinear}{colineales}
($\vec v = \frac32 \vec u$).

(b) $2 \times 10 - (-5) \times (-4) = 20 - 20 = 0$: \omterm{def:g10:vectors:collinear}{colineales}
($\vec v = -2\vec u$).

(c) $3 \times 3 - 1 \times 1 = 8 \neq 0$: no son \omterm{def:g10:vectors:collinear}{colineales}.
\end{solution}

\begin{solution}{exo:g10:vectors:5}
$\vect{AB}\,(4, -2)$ y $\vect{DC}\,(6 - x_D,\ 3 - y_D)$. La condición
$\vect{AB} = \vect{DC}$ da $6 - x_D = 4$ y $3 - y_D = -2$, luego
$D(2, 5)$.
\end{solution}

\begin{solution}{exo:g10:vectors:6}
\emph{1.} $\vect{AB}\,(3, 2)$ y $\vect{AC}\,(6, -2)$, luego
$\vect{AB} + \vect{AC} = (9, 0)$ y $M = (-2 + 9,\ 1 + 0) = (7, 1)$.

\emph{2.} Por la regla del paralelogramo, $\vect{AM} = \vect{AB} +
\vect{AC}$ significa que $ABMC$ es un paralelogramo ($M$ es el cuarto
vértice, opuesto a $A$). Se puede comprobar: $\vect{AB}\,(3,2)$ y
$\vect{CM}\,(7-4,\ 1-(-1)) = (3, 2)$ son iguales.
\end{solution}

\begin{solution}{exo:g10:vectors:7}
Primer trío: $\vect{AB}\,(3, 2)$, $\vect{AC}\,(9, 6)$;
$3 \times 6 - 2 \times 9 = 0$: alineados ($\vect{AC} = 3\vect{AB}$).

Segundo trío: $\vect{AB}\,(2, -1)$, $\vect{AC}\,(8, -4)$;
$2 \times (-4) - (-1) \times 8 = -8 + 8 = 0$: también alineados
($\vect{AC} = 4\vect{AB}$).
\end{solution}

\begin{solution}{exo:g10:vectors:8}
$\vect{AB}\,(4, 2)$ y $\vect{DC}\,(6-2,\ 1-(-1)) = (4, 2)$: los dos
\omterm{def:g10:vectors:vector}{vectores} son \emph{iguales}, así que $ABCD$ es un paralelogramo y, en
particular, $(AB) \parallel (CD)$. La propiedad de paralelogramo es la más
fuerte: implica el paralelismo (la colinealidad de los \omterm{def:g10:vectors:vector}{vectores}), pero no
al revés.
\end{solution}

\begin{solution}{exo:g10:vectors:9}
\emph{1.} $\vect{AB}\,(6, 3)$, luego $\vect{AI} = \frac23\vect{AB} =
(4, 2)$ e $I(4, 2)$.

\emph{2.} $\vect{CB}\,(6-2,\ 3-5) = (4, -2)$, luego
$\vect{CJ} = \frac23\vect{CB} = \left(\frac83, -\frac43\right)$ y
$J\left(2 + \frac83,\ 5 - \frac43\right) = \left(\frac{14}{3},
\frac{11}{3}\right)$.

\emph{3.} $\vect{IJ}\,\left(\frac{14}{3} - 4,\ \frac{11}{3} - 2\right) =
\left(\frac23, \frac53\right)$ y $\vect{AC}\,(2, 5)$. Criterio de
colinealidad:
\[
\frac23 \times 5 - \frac53 \times 2 = \frac{10}{3} - \frac{10}{3} = 0 :
\]
los \omterm{def:g10:vectors:vector}{vectores} son \omterm{def:g10:vectors:collinear}{colineales}, así que $(IJ)$ es paralela a $(AC)$. (Tanto
$I$ como $J$ están a un tercio del camino hacia $B$ desde su lado: el
pequeño triángulo $IBJ$ es una reducción de $ABC$.)
\end{solution}

\begin{solution}{exo:g10:vectors:10}
Los \omterm{prop:g10:coordgeom:midpoint}{puntos medios} son
\[
I\left(\tfrac{x_A + x_B}{2}, \tfrac{y_A + y_B}{2}\right), \quad
J\left(\tfrac{x_B + x_C}{2}, \tfrac{y_B + y_C}{2}\right), \quad
K\left(\tfrac{x_C + x_D}{2}, \tfrac{y_C + y_D}{2}\right), \quad
L\left(\tfrac{x_D + x_A}{2}, \tfrac{y_D + y_A}{2}\right).
\]
Primeras \omterm{def:g10:coordgeom:system}{coordenadas} de los dos \omterm{def:g10:vectors:vector}{vectores}:
\[
x_J - x_I = \frac{x_B + x_C}{2} - \frac{x_A + x_B}{2}
= \frac{x_C - x_A}{2},
\qquad
x_K - x_L = \frac{x_C + x_D}{2} - \frac{x_D + x_A}{2}
= \frac{x_C - x_A}{2},
\]
y el mismo cálculo vale para las segundas \omterm{def:g10:coordgeom:system}{coordenadas}. Por tanto,
$\vect{IJ} = \vect{LK}$ (los dos valen $\frac12\vect{AC}$) e $IJKL$ es un
paralelogramo, sea cual sea el aspecto del cuadrilátero $ABCD$.
\end{solution}

\begin{solution}{pb:g10:vectors:1}
\textbf{1.} $\vect{AB} + \vect{BC} + \vect{CD} = \vect{AD}$; \quad
$\vect{MA} - \vect{MB} = \vect{MA} + \vect{BM} = \vect{BA}$; \quad
$\vect{AB} + \vect{BC} + \vect{CD} + \vect{DA} = \vect{AA} = \vec 0$
(reordenando la tercera \omterm{def:g10:vectors:sum}{suma}).

\textbf{2.} $\vect{OB} - \vect{OA} = \vect{AO} + \vect{OB} = \vect{AB}$
por Chasles.

\textbf{3.} $I$ es el \omterm{prop:g10:coordgeom:midpoint}{punto medio} de $[AB]$ exactamente cuando
$\vect{IA} = -\vect{IB}$, es decir, $\vect{IA} + \vect{IB} = \vec 0$.
Entonces, para cualquier $O$:
$\vec 0 = \vect{IO} + \vect{OA} + \vect{IO} + \vect{OB}$, luego
$2\,\vect{OI} = \vect{OA} + \vect{OB}$; y recíprocamente.

\textbf{4.} $\vect{AB} = (4, 2)$, y $D$ cumple $\vect{DC} = \vect{AB}$:
$D = C - (4, 2) = (2, -2)$; el mismo $D$ que en el
\cref{pb:g10:coordgeom:1}, una línea más rápido.

\textbf{5.} La traslación de \omterm{def:g10:vectors:vector}{vector} $\vec u$, exactamente la
transformación que producen dos espejos paralelos o dos medios giros.
\omterm{def:g10:functions:function}{Imagen} de $(1, 2)$: $(4, 1)$.

\textbf{6.} Desarrollando cada $\vect{GX}$ como $\vect{GO} + \vect{OX}$:
$3\,\vect{GO} + \vect{OA} + \vect{OB} + \vect{OC} = \vec 0$, luego
$\vect{OG} = \frac13(\vect{OA} + \vect{OB} + \vect{OC})$: uno y solo un
punto así, cuyas \omterm{def:g10:coordgeom:system}{coordenadas} son las medias de las de los tres vértices.

\textbf{7.} $\vect{AB} + \vect{AC} =
(\vect{AG} + \vect{GB}) + (\vect{AG} + \vect{GC})
= 2\vect{AG} + (\vect{GB} + \vect{GC})
= 2\vect{AG} - \vect{GA} = 3\vect{AG}$, de donde
$\vect{AG} = \frac13(\vect{AB} + \vect{AC})$. La fórmula del \omterm{prop:g10:coordgeom:midpoint}{punto medio}
(pregunta 3, vista desde $A$) da
$\vect{AK} = \frac12(\vect{AB} + \vect{AC})$. Comparando:
$\vect{AG} = \frac23\vect{AK}$; $G$ está sobre la mediana $[AK]$, a dos
tercios del camino desde $A$.

\textbf{8.} La condición que define $G$ trata a $A$, $B$ y $C$ por igual,
así que el mismo cálculo desde $B$ o desde $C$ coloca el \emph{mismo} $G$
a dos tercios de esas medianas: la concurrencia sale gratis. La
demostración con el paralelogramo era ingeniosa, pero improvisada para el
caso; las \omterm{def:g10:coordgeom:system}{coordenadas} calculan, pero oscurecen; los \omterm{def:g10:vectors:vector}{vectores} hacen visible
la simetría y reducen la demostración a dos líneas.

\textbf{9.} $G\left(\frac{1+5+3}{3}, \frac{1+2+6}{3}\right) = (3, 3)$;
$K(4, 4)$; $\vect{AG} = (2, 2)$ y $\vect{AK} = (3, 3)$: determinante
$2 \times 3 - 2 \times 3 = 0$, \omterm{def:g10:vectors:collinear}{colineales}, y en efecto
$\vect{AG} = \frac23\vect{AK}$.

\textbf{10.} $G'\left(\frac{2 + 5 + 3}{4}, \frac{2 + 2 + 6}{4}\right) =
(2.5,\ 2.5)$: atraído desde $G(3,3)$ hacia el vértice duplicado $A(1,1)$,
como debe hacer un punto de equilibrio.

\textbf{11.} $3 \times 4 - (-2)(-6) = 12 - 12 = 0$: \omterm{def:g10:vectors:collinear}{colineales}
($\vec v = -2\vec u$). $2 \times 9 - 5 \times 4 = -2 \neq 0$: no
\omterm{def:g10:vectors:collinear}{colineales}.

\textbf{12.} $\vect{AB} = (2, 3)$, $\vect{AC} = (6, 9)$: determinante
$2 \times 9 - 3 \times 6 = 0$: alineados.

\textbf{13.} $\vect{OI} = \frac12(\vect{OA} + \vect{OB})$ y
$\vect{OK} = \frac12(\vect{OC} + \vect{OD})$, así que el \omterm{prop:g10:coordgeom:midpoint}{punto medio} de
$[IK]$ tiene \omterm{def:g10:vectors:vector}{vector} de posición
$\frac12(\vect{OI} + \vect{OK}) = \frac14(\vect{OA} + \vect{OB} +
\vect{OC} + \vect{OD})$; y el cálculo para $[JL]$ da exactamente el mismo
\omterm{def:g10:vectors:vector}{vector}. Las dos bimedianas (y, por Varignon, las diagonales del
paralelogramo $IJKL$) comparten ese centro: el baricentro del
cuadrilátero.

\textbf{14.} $\vect{MN} = \vect{AN} - \vect{AM} =
\frac13\vect{AC} - \frac13\vect{AB} = \frac13\vect{BC}$: luego
$(MN) \parallel (BC)$ y $MN = \frac13 BC$.

\textbf{15.} Con $\vect{AM} = k\,\vect{AB}$ y $\vect{AN} = k\,\vect{AC}$:
$\vect{MN} = k\,\vect{BC}$; una sola línea que contiene el teorema de la
base media ($k = \frac12$) y el de Tales (cualquier $k$): la razón
atraviesa la resta sin inmutarse.

\textbf{16.} Resultante $(3, 4)$, módulo $\sqrt{9 + 16} = 5$ km/h: la
barca avanza realmente hacia el nordeste, tirando al norte, siguiendo la
hipotenusa del triángulo $3$--$4$--$5$.

\textbf{17.} $\vec u + \vec v + \vec w = (0, 0)$: la anilla no se mueve;
está en equilibrio. A lo largo de un polígono cerrado
$A_1 A_2 \dots A_n A_1$, Chasles condensa la \omterm{def:g10:vectors:sum}{suma} en
$\vect{A_1 A_1} = \vec 0$: dar la vuelta te devuelve al punto de partida.

\textbf{18.} $\vec w = -(\vec u + \vec v) = -(3, 2) = (-3, -2)$.

\textbf{19.} Resultante $(1.5,\ 2)$, módulo $\sqrt{2.25 + 4} = 2.5$ m/s.
Travesía: $50$ m a $2$ m/s en la dirección transversal: $25$ s; deriva:
$1.5 \times 25 = 37.5$ m río abajo.

\textbf{20.} Demostrados de nuevo este fin de semana: el teorema de los
dos tercios del baricentro (pregunta 7), el centro de las bimedianas y de
Varignon (pregunta 13) y Tales con la base media (preguntas 14 y 15). Lo
que las flechas todavía no saben hacer: medir \emph{ángulos} (y, con
ellos, longitudes en direcciones arbitrarias); el producto escalar del
curso que viene es justo el instrumento que falta.
\end{solution}
